Turunan parsial dari

\mavergleichskettedisp
{\vergleichskette
{ \psi \begin{pmatrix}  a  \\b   \end{pmatrix}
}
{ =} { { \frac{   1 }{  -2a+1+a^2+b^2 } } \begin{pmatrix}  -2b \\ 1-a^2-b^2   \end{pmatrix}
}
{ } {
}
{ } {
}
{ } {
}
}
{}{}{}
adalah

\mavergleichskettealign
{\vergleichskettealign
{ { \frac{   \partial   \psi  }{  \partial  a  } }
}
{ =} { { \frac{   1 }{  \left(  -2a+1+a^2+b^2  \right)^2  } } \begin{pmatrix}  2b \left(  2a-2  \right)  \\ -2a \left(  -2a+1+a^2+b^2  \right) - \left(  1-a^2-b^2  \right) \left(  -2+2a  \right)    \end{pmatrix}
}
{ =} { { \frac{   1 }{  \left(  -2a+1+a^2+b^2  \right)^2  } } \begin{pmatrix}  2b \left(  2a-2  \right)  \\ 4a^2 -4a+2 \left(  1-a^2-b^2  \right)    \end{pmatrix}
}
{ =} { { \frac{   1 }{  \left(  -2a+1+a^2+b^2  \right)^2  } } \begin{pmatrix}  4ab -4b  \\ 2a^2 -4a+2 -2 b^2     \end{pmatrix}
}
{ } {
}
}
{}
{}{}
dan

\mavergleichskettealign
{\vergleichskettealign
{ { \frac{   \partial   \psi  }{  \partial  b  } }
}
{ =} { { \frac{   1 }{  \left(  -2a+1+a^2+b^2  \right)^2  } } \begin{pmatrix}  -2 \left(  -2a+1+a^2+b^2  \right) +4b^2   \\ -2b \left(  -2a+1+a^2+b^2  \right) - 2b \left(  1-a^2-b^2  \right)     \end{pmatrix}
}
{ =} { { \frac{   1 }{  \left(  -2a+1+a^2+b^2  \right)^2  } } \begin{pmatrix}  4a  -2-2a^2+2b^2   \\ 4ab -4b    \end{pmatrix}
}
{ } {
}
{ } {
}
}
{}
{}{.}
